\section{Linear Markov Setting}
\label{sec:linear-markov-setting}

\subsection{Stochastic and signed retractions}

Fix a finite state set \(\{1,\ldots,n\}\).  Vectors are row vectors.
The all-ones column is \(\one\), the identity matrix is \(I_n\), and
\(J=\one\one^{\mathsf T}\) is the all-ones matrix; source-note shorthand
\(2Q-1\) below means \(2Q-I_n\), not \(2Q-J\).  The simplex
\[
  \DeltaN=\{x\in \R^n:\ x_k\ge 0\ \hbox{for all }k,\quad x\one=1\}
\]
is the standard simplex.  A matrix \(Q\in M_n(\R)\) is
\emph{row-stochastic} if every row lies in \(\DeltaN\), equivalently
\(Q\one=\one\) and \(Q_{ik}\ge 0\) for all \(i,k\).  A
\emph{stochastic idempotent} is a row-stochastic matrix \(S\) with
\(S^2=S\).  The original almost-idempotent defect is
\[
  \eta(Q)=\lVert Q^2-Q\rVert_{\infty\to\infty},
\]
where \(\lVert A\rVert_{\infty\to\infty}\) is the maximum absolute row
sum norm.

The exact model used in this report is a \emph{signed affine
retraction}: a real matrix \(P\) satisfying
\[
  P\one=\one,\qquad P^2=P,
\]
with no entrywise positivity assumption.  Its \(i\)th row is denoted
\(p_i\).  We use the same symbol \(p_i\) both for the matrix row and
for the point of the affine hyperplane \(\{x\in \R^n:x\one=1\}\) that
it defines.  The \emph{row polytope} is
\[
  K=\conv\{p_1,\ldots,p_n\}.
\]
Thus \(K\) is a convex hull of rows of \(P\), not the simplex unless
the rows themselves are the coordinate vertices.

The identity \(P^2=P\) is the main geometric rule.  Row \(i\) of this
identity says
\[
  p_i=\sum_{k=1}^n P_{ik}\,p_k,\qquad \sum_{k=1}^n P_{ik}=1.
\]
So each row is an affine combination of all rows, with coefficients
given by the same row of \(P\).  If all coefficients in row \(i\) are
nonnegative, the identity only says \(p_i\in K\).  If some coefficients
are negative, the same identity is a signed balance law.  For example,
the coefficient pattern \(P_{i1}=1+\varepsilon\),
\(P_{i2}=-\varepsilon\), and all other \(P_{ik}=0\) reads
\[
  p_i=p_1+\varepsilon(p_1-p_2),
\]
so the row lies on the affine line through \(p_2,p_1\), beyond \(p_1\).
This is the elementary geometry behind the negative-mass parameter.

\subsection{The L1 geometry and negative mass}

Distances in the report are measured in the \(\ellone\) norm,
\[
  \lVert x\rVert_1=\sum_{k=1}^n |x_k|.
\]
For a set \(C\subset\R^n\), write
\[
  \distone(x,C)=\inf_{y\in C}\lVert x-y\rVert_1,
\]
and set
\[
  D_K=\diamone(K)=\sup_{x,y\in K}\lVert x-y\rVert_1.
\]
The dual unit ball for this metric is the \(\ell^\infty\) ball; later
separators are normalized using this duality.

For a scalar \(a\), write
\[
  a_+=\max(a,0),\qquad a_-=\max(-a,0).
\]
Thus \(a=a_+-a_-\) and \(|a|=a_++a_-\).  The \emph{row negative mass}
of \(P\) at row \(i\) is
\[
  \nu_i=\sum_{k=1}^n (P_{ik})_-,
\]
and the global max-negative parameter is
\[
  \delta=\max_i \nu_i,\qquad \tau=\sqrt{\delta}.
\]
Because \(P\one=\one\), the positive coefficient mass in row \(i\) is
\[
  \sum_k (P_{ik})_+=1+\nu_i.
\]
The phrase \emph{max-neg units} means that all estimates are expressed
using this single parameter \(\delta\).  Unless a statement explicitly
says that a constant is construction-specific or finite-corner-specific,
all constants are universal: independent of \(n\), the rank of \(P\),
and the number of rows realizing the same point.

A coordinate package that appears in the computational notes is a
factorization
\[
  P=\Lambda\mathsf R,\qquad \mathsf R\Lambda=I.
\]
Here \(\mathsf R\) is a frame or realization matrix and \(\Lambda\) is
a coefficient or load matrix.  This factorization is only a way to
parametrize exact retractions.  The invariant objects in the report are
the rows \(p_i\), their \(\ellone\) geometry, and the signed row
coefficients of \(P\).  When the numerical sections quote a realized
instance, \(\delta\) and \(\tau\) are recomputed from that exact \(P\);
they are not inherited from a design parameter.

\subsection{Two imported reductions}

\begin{proposition}[Imported signed/stochastic equivalence]
\provedbadge
\label{prop:classical-equiv-consumed}
\texttt{lem-classical-equiv} is imported as a black-box theorem.  Its
source contract, with all symbols expanded, is:
\begin{quote}\small
The signed-idempotent and stochastic-idempotent formulations of classical
stability are equivalent up to universal constants.  For sufficiently
small \(\eta\), a row-stochastic \(Q\) with
\(\lVert Q^2-Q\rVert_{\infty\to\infty}\leq\eta\) gives
\(P=\theta(2Q-I_n)\), where
\(\theta(X)=\tfrac12(I_n+\operatorname{sgn}X)\) and
\(\operatorname{sgn}X=X(X^2)^{-1/2}\) is the finite-dimensional
spectral/sign transform used in the equivalence.  This \(P\) is a signed affine
retraction with
\(\lVert P-Q\rVert_{\infty\to\infty}\leq C\eta\) and negative mass
\(\delta\leq C\eta\).  Conversely, if \(P\one=\one\), \(P^2=P\), and
\(\delta=\max_i\sum_k(P_{ik})_-\), define
\(p_i^+=((P_{ik})_+)_{k=1}^n\), \(a_i=\sum_k(P_{ik})_-\), and
\(q_i=p_i^+/(1+a_i)\).
The row-normalized positive-part matrix \(Q\), with rows \(q_i\), is
row-stochastic and satisfies
\(\lVert P-Q\rVert_{\infty\to\infty}\leq 2\delta\) and
\(\lVert Q^2-Q\rVert_{\infty\to\infty}\leq 6\delta+4\delta^2\).
\end{quote}
All constants in this import are universal.  Thus the original
dimension-free problem for row-stochastic \(Q\) with small
\(\eta(Q)=\lVert Q^2-Q\rVert_{\infty\to\infty}\) may be studied through
exact signed affine retractions \(P\) with small \(\delta=\max_i\nu_i\).
\end{proposition}

\begin{proposition}[Imported cluster theorem]
\provedbadge
\label{prop:cluster-consumed}
\texttt{thm-cluster} is imported as a black-box theorem.  Its source
contract, with the local symbols made explicit, is:
\begin{quote}\small
There are universal \(\delta_0,C>0\) such that a signed affine
retraction \(P\) with negative mass at most
\(\delta\leq\delta_0\) has the following clustering consequence.  Suppose
there are representative row vertices
\(r^1=p_{i_1},\ldots,r^m=p_{i_m}\), pairwise
\(2\rho\)-separated in \(\ell^1\), and each representative is
\((\rho,\kappa)\)-exposed.  Let \(U_a\) be the associated disjoint
clusters of rows near \(r^a\), and suppose every row outside the clusters
lies within \(\gamma\) of \(\conv\{r^1,\ldots,r^m\}\).  Then there is a
stochastic idempotent \(E\) with
\(\lVert P-E\rVert_{\infty\to\infty}\leq
C(\rho+\gamma+\delta/\kappa+\delta)\); so
\(\rho,\gamma=O(\sqrt\delta)\), \(\kappa\geq c\sqrt\delta\) give
\(O(\sqrt\delta)\), \(C\) independent of \(m\), \(n\), and the number of
transient rows.
\end{quote}
This report does not reprove clustering; later sections only use the
contract above to pass from exposed-hull control to the original
classical stability target.
\end{proposition}

The consequence is methodological but important.  From this point on,
the report studies exact matrices \(P\) satisfying \(P\one=\one\) and
\(P^2=P\), with small \(\delta\).  Almost-idempotent matrices \(Q\) only
enter through \cref{prop:classical-equiv-consumed}; stochastic
idempotents are the target objects to which the signed rows must be
rounded or clustered.

\subsection{Where the remaining branch starts}

The next sections define exposed row vertices, the well-exposed set
\(\mathcal W\), its convex hull \(C_W=\conv(\mathcal W)\), and the
\emph{hidden height} \(H=\max_i\distone(p_i,C_W)\).  They also define,
for a height-maximizing hidden row vertex \(v\), the geometric mass
\(\widetilde\sigma_v\): positive \(P_{vk}\)-mass carried by rows lying
outside the exposed hull, with the small halo convention stated there.
Those are not extra data beyond the signed retraction.  They are
measurements made from the same row identity
\[
  p_v=\sum_k P_{vk}p_k.
\]

\begin{openproblem}[Residual branch]
\openbadge
\label{op:residual-branch-preview}
After the exposed geometry is introduced, the final open form is the
height-localization claim: prove a universal \(a>0\) with
\(\delta\ge aH^2\).  Equivalently in the post-design formulation, it
remains to exclude hidden top vertices with
\(\widetilde\sigma_v>\tau\) at height \(H>B\tau\), for a universal
constant \(B\).  The complementary branch
\(\widetilde\sigma_v\le\tau\) is no longer open:
\(\provedbadge\) the height cap gives \(H\le 3\tau\) for
\(\delta\le 1/4\), hence \(\delta\ge H^2/9\) directly.  The exclusion
must be height-conditional:
\(\modauditbadge\) exact rational low-height all-shallow optimal faces
exist with \(\widetilde\sigma_v>0\), but only at \(H=O(\delta)\).
\end{openproblem}

This is why the basic setting is deliberately linear.  The remaining
obstruction is not a failure of \(P^2=P\), nor a failure of the Markov
normalization \(P\one=\one\).  It is a signed quantitative
Baake--Sumner stability question for affine row identities near the
exposed hull.  \(\provedbadge\) The pushed-witness support-cleanup route
is already recorded as unable to close the remaining branch, because
negative leakage breaks the dual optimality needed by that route.
